\section{卷积网络的结构与局部计算}
\label{l11:sec:structure}

卷积神经网络（Convolutional Neural Network，CNN）用局部加权求和提取特征，再通过非线性变换和空间下采样构成多层计算。实现这类网络时，需要同时确定三个相互联系的对象：每个输出值依赖哪些输入，张量中的元素如何存放，以及哪些输出可以交给不同线程独立计算。卷积层的数学定义决定数据依赖，而线程映射和存储组织决定这些依赖如何在图形处理器（Graphics Processing Unit，GPU）上得到高效执行。

\subsection{从全连接计算到局部连接}

多层感知机（Multi-Layer Perceptron，MLP）的全连接层把输入写成向量 $\symbf{x}\in\mathbb{R}^n$。若该层有 $m$ 个输出，则第 $j$ 个输出为 $a_j=\phi(\sum_{i=0}^{n-1}V_{j,i}x_i+\beta_j)$，其中 $V\in\mathbb{R}^{m\times n}$ 是权重矩阵，$\beta_j$ 是偏置，$\phi$ 是逐元素激活函数。一个输出连接全部 $n$ 个输入，因此该层有 $mn$ 个权重和 $m$ 个偏置。

例如，把一幅 $28\times28$ 的灰度图像展平为 $784$ 维向量，先连接到 $10$ 个隐藏神经元，再连接到 $10$ 个输出神经元，总参数量为 $(784\times10+10)+(10\times10+10)=7\,960$。当图像增大到 $250\times250$，输入维数为 $62\,500$；若隐藏层也有 $62\,500$ 个神经元，仅第一层就需要 $62\,500^2=3\,906\,250\,000$ 个权重。按每个权重占 $4$ 字节计算，权重存储约为 $15.625$ GB，还没有计入偏置、激活值以及训练时的其他数组。

图像中的局部边缘、纹理等模式通常可以通过一个较小区域内的像素关系描述。卷积层因此让每个输出只连接输入中的局部窗口，形成稀疏交互（sparse interaction）。一个 $K\times K$ 的窗口在单个输入通道上只需要 $K^2$ 次对应元素相乘，而不需要访问整幅图像。这个窗口是该输出相对于上一层的感受野（receptive field）。

局部连接与多层组合可以同时存在。在窗口完整落入有效输入范围的位置，若每层都使用步幅为 $1$ 的 $3\times3$ 卷积，第一层输出依赖输入中的 $3\times3$ 区域；第二层把第一层相邻的 $3\times3$ 个位置组合起来，这些位置对应原输入中的 $5\times5$ 区域。继续堆叠 $\ell$ 层时，沿任一空间轴的感受野宽度为 $r_\ell=1+\ell(K-1)$。这个关系假设各层核宽均为 $K$、没有空洞且步幅均为 $1$；它说明局部运算可以逐层建立较大范围的数据依赖。

\subsection{参数共享与平移等变性}

卷积层还把同一组权重应用到不同空间位置，这称为参数共享（parameter sharing）。对于固定的输入通道和输出通道，窗口向右或向下移动时，卷积核系数保持不变。输出位置改变的是取用的输入数据，而不改变模型参数。若输入有 $C$ 个通道、输出有 $M$ 个通道、每个二维核为 $K\times K$，则权重总数为 $MCK^2$；每个输出通道再有一个偏置时，参数总数为 $MCK^2+M$。这里没有乘输出的空间位置数。

例如，单通道图像经过 $16$ 个 $5\times5$ 核产生 $16$ 张特征图，只需 $16\times1\times25=400$ 个权重，含偏置时为 $416$ 个参数。只要输入仍有一个通道，这组权重就能作用于不同高度和宽度的图像，输出特征图的尺寸随之改变。若后续把全部空间位置展平再送入一个固定大小的全连接层，则该全连接层的输入维数仍然固定；因此，卷积核可复用于不同图像尺寸，与整个网络能否接收任意尺寸，是两个需要分别检查的条件。

参数共享还产生平移等变性（translation equivariance）：输入中的局部模式平移后，对应响应的位置随之平移。以单通道、步幅为 $1$ 的二维卷积为例，记其算子为 $\mathcal{C}_F$，令 $\Delta_h,\Delta_w$ 为整数。二维平移算子满足 $(T_{\Delta}X)[h,w]=X[h-\Delta_h,w-\Delta_w]$，卷积核不随位置改变，且所比较位置的取值和边界处理均有效，则
\begin{equation}
  \mathcal{C}_F(T_{\Delta}X)[h,w]
  =\sum_{p,q}F[p,q]X[h+p-\Delta_h,w+q-\Delta_w]
  =(T_{\Delta}\mathcal{C}_F(X))[h,w].
  \label{l11:eq:equivariance}
\end{equation}
左右两边使用完全相同的乘积项，只是描述坐标不同。等变性保留特征出现的位置；如果需要输入平移后最终分类结果保持不变，还需要后续的汇聚和分类过程。有限图像的裁剪、填充方式，以及步幅大于 $1$ 的下采样，会限制这种逐像素平移对应关系。

\subsection{空间维度、通道与特征图}

空间维度描述卷积窗口移动的方向，通道（channel）描述同一位置上并列的不同信号。一幅彩色图像有高度和宽度两个空间轴，在每个位置有红、绿、蓝三个颜色通道。三通道二维卷积的窗口沿两个空间轴移动，并在三个通道之间求和；它不会因为存在三个通道就变成沿三个空间轴移动的三维卷积。

\begin{table}[htbp]
\centering
\caption{空间轴与通道的不同含义}
\label{l11:tab:dimensions}
\begin{tabularx}{\linewidth}{L{0.17\linewidth}L{0.30\linewidth}X}
\toprule
数据对象 & 卷积移动的轴 & 通道的含义\\
\midrule
音频波形 & 一维时间轴 & 单路或多路音频信号\\
骨架动作序列 & 一维时间轴 & 不同关节的角度或坐标分量\\
彩色图像 & 二维高度、宽度 & 同一像素的红、绿、蓝分量\\
时频表示 & 二维时间、频率 & 同一时频位置的不同测量分量\\
体数据 & 三维空间坐标 & 同一体素的一个或多个测量分量\\
彩色视频 & 二维空间与一维时间 & 同一时空位置的颜色分量\\
\bottomrule
\end{tabularx}
\end{table}

输出通道也称输出特征图（feature map）。每个输出通道有一整组覆盖所有输入通道的权重：对各输入通道分别计算局部加权和，再把这些部分结果相加，得到一个输出值。因而 $C$ 个输入通道与 $M$ 个输出通道之间共有 $C\times M$ 个二维卷积核，但最终输出是 $M$ 张二维特征图。不同的输出通道拥有不同权重，可以提取不同模式；同一输出通道的权重在所有图像和空间位置上共享。

\begin{figure}[htbp]
\centering
\begin{tikzpicture}[>=Latex,font=\small,node distance=9mm]
  \foreach \i/\y in {0/1.2,1/0,2/-1.2}{
    \node[draw=UIUCBlue,minimum width=2.2cm,minimum height=.72cm] (x\i) at (0,\y) {$X_{b,\i}$ 的局部窗口};
    \node[draw=UIUCBlue,minimum width=2.6cm,minimum height=.72cm] (f\i) at (4.0,\y) {与 $F_{m,\i}$ 相乘并求和};
    \draw[->] (x\i) -- (f\i);
  }
  \node[draw=UIUCOrange,circle,minimum size=.8cm] (sum) at (6.7,0) {$+$};
  \node[draw=UIUCBlue,minimum width=1.75cm,minimum height=.8cm] (y) at (9.0,0) {$Y_{b,m,h,w}$};
  \foreach \i in {0,1,2}{\draw[->] (f\i.east) -- (sum.west);}
  \draw[->] (sum) -- (y);
\end{tikzpicture}
\caption{三个输入通道的贡献相加，生成固定输出通道中的一个空间位置。对每个新的输出通道，都需要另一组覆盖三个输入通道的卷积核。}
\label{l11:fig:channels}
\end{figure}

\subsection{卷积、池化与分类层的组合}

卷积层产生包含空间位置信息的特征图，激活函数对其中每个标量执行非线性变换，池化（pooling）则把局部区域汇聚为较少的输出值。多次组合这些操作后，网络可以将较大感受野中的信息交给全连接层或其他分类结构。每个模块都接收和产生张量，上一层输出的通道数与空间尺寸直接决定下一层计算的输入形状。

经典卷积网络（LeNet-5）的一个形状序列是：$32\times32$ 单通道图像经过卷积得到 $6\times28\times28$，下采样为 $6\times14\times14$，再卷积为 $16\times10\times10$，再下采样为 $16\times5\times5$，之后形成 $120$ 维、$84$ 维的表示并产生 $10$ 类输出。空间尺寸的收缩来自有效卷积和下采样，通道数的变化来自输出卷积核组数。较深的网络可以采用更多卷积层，也可以引入跨层连接或上采样分支；这些网络结构改变了张量之间的连接方式，卷积层内部仍可从同一套局部求和与通道组合规则出发实现。

卷积的局部连接主要减少每个输出所需的输入范围，参数共享主要减少独立权重数量，池化主要减少空间位置数。这三种机制作用于不同对象，因此需要分别分析参数量、计算量和中间张量存储量。
